\documentclass[11pt,a4paper]{article}
\usepackage[a4paper,margin=25mm]{geometry}
\usepackage[T1]{fontenc}
\usepackage{lmodern}
\usepackage{microtype}
\usepackage[hidelinks]{hyperref}
\setlength{\parindent}{0pt}
\setlength{\parskip}{0.75em}
\begin{document}

\begin{flushright}
Angran Xia\\
School of Chemistry and Chemical Engineering\\
Shaanxi Normal University\\
Xi'an, Shaanxi 710119, China\\
\href{mailto:haangyeon@snnu.edu.cn}{haangyeon@snnu.edu.cn}\\[1em]
\today
\end{flushright}

Dear Editors of \emph{Genes},

Please consider our manuscript, ``Cross-model uncertainty-aware minimal editing of
cis-regulatory elements for cell-type-selective design,'' for publication as a Research
Article. The study addresses a biological problem in functional genomics: how to
retune natural cis-regulatory elements (CREs) with a small, interpretable number of
nucleotide substitutions while preserving a cell-type-selective regulatory signal.

Existing sequence-design workflows have made sequence-to-activity prediction highly
effective, but commonly select edits with the same predictor that defines the objective
or return a single unconstrained high-scoring sequence. SafeEdit-CRE introduces a
procedurally separated framework for minimal editing: candidate generation uses one
predictor, cross-model review uses an architecture-diverse ensemble, and final
evaluation uses a held-out multi-kernel model. It combines the published
cell-type-specificity predictor with validation-calibrated uncertainty,
sequence-domain controls, and constrained beam search on natural 200-nt CRE scaffolds,
prioritizing variants with consistent predictions across separately trained,
architecture-diverse models.

The principal findings are:
\begin{itemize}
  \item Across 600 held-out natural CRE parents in K562, HepG2, and SK-N-SH cells,
  SafeEdit-CRE increases cross-model specificity-margin gain from 0.822 for greedy
  editing to 0.877 (paired difference 0.055; 95\% CI 0.040--0.071).
  \item Against a compute-matched primary-model beam, predicted specificity remains
  closely matching (difference $-0.004$; 95\% CI $-0.011$--$0.004$), while cross-model
  uncertainty decreases from 0.335 to 0.314.
  \item In a locked, non-overlapping 90-parent design-quality benchmark, the fraction
  passing all prespecified sequence and model-agreement checks increases from 38.8\%
  for greedy editing to 60.6\% for SafeEdit-CRE, yielding 74 Tier A variants
  prioritized for reporter screening.
\end{itemize}

These results make a focused contribution to the scope of \emph{Genes}: they connect
sequence modeling with cis-regulatory grammar, cell-type-specific transcriptional
control, and reproducible functional-genomics design. The public code, checkpoints,
candidate library, and frozen summary inputs are available at
\href{https://github.com/Haangyeon/SafeEdit-CRE}{GitHub} and in the versioned Zenodo
archive (DOI: \href{https://doi.org/10.5281/zenodo.21458489}{10.5281/zenodo.21458489}).

The manuscript is original, is not under consideration elsewhere, and has been
approved by the author. The author declares no competing interests. Thank you for
considering this work.

Sincerely,\\[0.8em]
Angran Xia

\end{document}
